%% Bernoulli - Ketten
\subsection{Bernoulli-Ketten}\index{Verteilung!Bernoulli}\index{Bernoulli}
(Binomialverteilung)

\TRAINER{Einstiegsvideo: ``EinfachMathe!''\\
  \texttt{https://www.youtube.com/watch?v=WC5O317JzW0}}

\subsubsection{Einstiegsaufgabe}

\paragraph{Genau zwei Sechser} Ein Würfel wird \textbf{dreimal} hintereinander geworfen.

Wie groß ist die Wahrscheinlichkeit, dass genau zwei Sechser dabei sind?


Zeichnen Sie das zugehörige Baumdiagramm:

\TNT{18}{
\bbwCenterGraphic{15cm}{geso/stoch/img/genau2sechser.png}


  Lösung mit Baum oder  Bernoulli-Formel (Binomialverteilung).

  Genau zwei Sechser: $\left(\frac16\right)^2\cdot{}\frac56 +
  \left(\frac16\right)^2\cdot{}\frac56 +
  \left(\frac16\right)^2\cdot{}\frac56 =3\cdot{}
  \left(\frac16\right)^2\cdot{}\frac56= \frac{5}{72}\approx 0.06944$%%
%%
  \vspace{10cm}%%
}%% END TNT
\newpage


\subsubsection{Bernoulli-Experiment}\index{Experiment!Bernoulli}\index{Bernoulli-Experiment}

\bbwCenterGraphic{5cm}{geso/stoch/img/BernoulliJakobI_phi.jpg}
\begin{center}\textit{Jakob Bernoulli I (1654-1705)}\end{center}

Obiges Würfelexperiment ist ein typisches
Bernoulli-Experiment\footnote{Jakob I Bernoulli, Schweizer Mathematiker,
  1654-1705}, d.\,h.:

\begin{definition}{Bernoulli-Experiment}{}
\begin{itemize}
\item Das Experiment wird $n$-mal durchgeführt\footnote{Daher wird dieser
Bernoulli-Prozess manchmal auch Bernoulli-Kette genannt. Quelle
\texttt{www.wikipedia.org} 2021-04-08}.
\item Das Einzelexperiment hat genau zwei Ausgänge: Erfolg/Misserfolg.
\item Jedes Einzelexperiment hat für die erfolgreichen Ausgänge die gleiche
      Wahrscheinlichkeit $p$ (somit ist die Wahrscheinlichkeit für den
      Misserfolg jeweils gleich $1-p$).
\item Die Einzelexperimente sind voneinander unabhängig.
\end{itemize}
\end{definition}

\begin{bemerkung}{Bernoulli-Experiment}{}
Typischerweise sind wir bei Bernoulli-Experimenten an der
Zufallsvariable $X$ interessiert, welche angibt, wie oft ein Erfolg
eingetreten ist: $X$ hat also die Werte 0, 1, 2, 3, 4, ..., $n$.
\end{bemerkung}

\textbf{Kein} Bernoulli-Experiment ist beispielsweise das Ziehen von zwei Bällen aus
einer Urne mit drei gelben und sechs orangefarbenen Bällen \textbf{ohne} diese
jeweils zurückzulegen; denn dabei verändern sich die
Wahrscheinlichkeiten der verbleibenden Bälle nach jedem Herausnehmen.
\newpage


\paragraph{Beispiele von Bernoulli-Experimenten}

\begin{itemize}
\item
In einer Urne liegen zehn Kugeln: vier blaue und sechs grüne. Ein
Einzelexperiment besteht darin, genau eine Kugel zu ziehen, die Farbe
zu notieren und die Kugel sogleich wieder zurückzulegen.
Das Gesamtexperiment besteht darin, fünf Einzelexperimente
durchzuführen. Jedesmal ist die Wahrscheinlichkeit, eine blaue Kugel
zu erwischen, gleich groß und wir sprechen von einem Bernoulli-Experiment.

\item Vierfaches Drehen am Glücksrad: Wie groß ist die
  Wahrscheinlichkeit, genau zweimal «Gewinn» zu erzielen, wenn pro
  Drehung ein Gewinn mit 28 \% Wahrscheinlichkeit auftritt?

\item Siebenfaches Würfeln mit einem Spielwürfel, bei dem wir
  interessiert sind, wie groß die Wahrscheinlichkeit ist, dass genau fünf mal ein kleinerer Wurf als eine \epsdice{3} gewürfelt wird.
  
\item
Es müssen aber nicht immer Würfel, Münzen oder Kugeln in Urnen sein:
Der Torwart Alex Calderoni\footnote{Alex Calderoni * 1976: Bekannter
  italienischer Torwart beim Fußballspiel (Quelle Wikipedia 2022).}
fängt womöglich einen Ball mit einer
Wahrscheinlichkeit von 38\%.
Wie groß ist die Wahrscheinlichkeit, dass er genau sechs von zehn Torschüssen
abfangen kann?
\end{itemize}

\TRAINER{Diese Aufgaben allenfalls auch aufs Stochastik-Aufgabenblatt.}
\newpage


\subsubsection{Formel von Bernoulli}\index{Bernoulli!Formel von}\index{Formel von Bernoulli}

\TRAINER{Einstiegs-Viedo: \texttt{https://www.youtube.com/watch?v=WC5O317JzW0}}


\begin{gesetz}{Bernoulli-Formel (Binomialverteilung)}{}
  $$P(X=k) = {{n}\choose {k}}\cdot{}p^k\cdot{}(1-p)^{n-k}$$


%%\renewcommand{\arraystretch}{2}
\begin{bbwFillInTabular}{rcl}
  $n$ &=& \TRAINER{Anzahl Durchführungen (Baumtiefe)}\\
  $p$ &=& \TRAINER{Gewinnwahrscheinlichkeit pro Durchführung}\\
  $k$ &=& \TRAINER{Anzahl geforderte Gewinne}\\
\end{bbwFillInTabular}
\end{gesetz}


Bei der sog. Bernoulli-Kette handelt es sich um ein
Bernoulli-Experiment
in einem mehrstufigen Baum ($n$ Versuche). Dabei interessiert
uns, wie groß die Wahrscheinlichkeit ist, dass von den beiden
möglichen Ausgängen (Treffer / Niete) einer davon \textbf{genau} $k$-mal auftritt.

\begin{bemerkung}{Binomialkoeffizient}{}
  Wenn ich den Baum ablaufe, gibt es immer $n$ Verzweigungen.

  An genau $k$
Stellen will ich Richtung «Treffer» abzweigen. Wie viele Varianten
habe ich, $k$ Stellen aus $n$ auszuwählen?

\vspace{4mm}

Es gibt \LoesungsRaumLen{40mm}{$n \choose k$} Möglichkeiten, $k$ Verzweigungen aus
$n$ auszuwählen.
\end{bemerkung}

\newpage
\begin{beispiel}{Sieben mal Würfeln}{}
Wie groß ist die Wahrscheinlichkeit, mit einem Würfel bei
siebenmaligem Werfen genau fünf mal eine der Zahlen \epsdice{1} oder
\epsdice{2} zu erreichen?
\end{beispiel}
Dabei legen wir fest (vgl. Beispiel oben: Siebenmaliges Würfeln):
\begin{itemize}

\item
  $n$ = Anzahl mögliche Durchführungen; hier 7

\item
  $p$ = Wahrscheinlichkeit, dass eine \epsdice{1} oder \epsdice{2}
  geworfen wird. Hier $p = \frac26$.\\
  Somit ist die Misserfolgsquote $= (1-p)$; hier $1-\frac26=\frac46$.


\item
  $k$ = Anzahl der gewünschten ``Treffer''; hier 5, denn wir wollen fünf
Mal eine \epsdice{1} oder eine \epsdice{2} erzielen.
\end{itemize}

Nun gilt

In obigem Beispiel ergibt sich:
\vspace{2mm}
$$P(X=5) = \LoesungsRaumLen{70mm}{{7\choose 5} \cdot{}
  \left(\frac26\right)^5 \cdot{} \left(1-\frac26\right)^{7-5}} =
\LoesungsRaumLen{60mm}{{7\choose 5} \cdot{} \left(\frac26\right)^5
  \cdot{} \left(\frac46\right)^2}$$


\textbf{Taschenrechner}:

\leserluft

Dies kann mit dem Taschenrechner gelöst werden:
Entweder direkt ...

$$\LoesungsRaumLang{( 7 \text{\,\,nCr\,\,} 5 ) *\left(\frac26\right)^5 *
\left(1-\frac26\right)^{7-5}} \approx \LoesungsRaum{0.03841 = 3.841\% =\frac{28}{729}}$$

... oder aber mit der eingespeicherten Formel bei
\tiprobutton{data_stat-reg-distr} und dort unter \texttt{DISTR: 4
  Binomialpdf: SINGLE}

... dann ...

\begin{tabular}{c@{ $:$ }l}
  $n=7$   & ... bei \texttt{n} \\
  $p=2/6$ & bei \texttt{p(\text{success})} und \\
  $k=5$   & bei \texttt{x} eintragen.
\end{tabular} 
\newpage


\subsection*{Aufgaben}

\aufgabenFarbe{Man würfelt 4-mal hintereinander mit einem gewöhnlichen Spielwürfel. Wie
gross ist die Wahrscheinlichkeit, dass man genau 2 Sechser gewürfelt
hat?}
\TNT{4}{
  $$P(X=2) = {4 \choose 2} \cdot{} \left(\frac{1}{6}\right)^2 \cdot{}
  \left(\frac{5}{6}\right)^{4-2} \approx 11.57\%$$
}%% end TNT


\aufgabenFarbe{Maturaprüfungen:\\2016 (GESO) Aufgabe 11 (Geburten)\\2017 (GESO) Aufgabe 17. (borkenkäferresistent)\\2018 (GESO) Serie 1 Aufg. 16. (Papayasendung)\\2018 (GESO) Serie 2 Aufg. 15. (Sprössling)\\2018 (GESO) Serie 3 Aufg. 14. (Ikosaeder)}


\newpage
\subsubsection{Kumulierte Wahrscheinlichkeit}

\subsubsection*{Motivation}

\begin{beispiel}{Hygienemasken}{}
  0.2\,\% der hergestellten Hygienemasken sind defekt. Ich kaufe 1000
  Stk. Wie groß ist die Wahrscheinlichkeit, dass weniger als 6 defekt sind?
  \end{beispiel}
\newpage

Basketballspieler «Basil Ballisti» trifft mit einer
Wahrscheinlichkeit von 85\%.

Die Wahrscheinlichkeit, dass er von 20 Würfen \textbf{genau} 17 Mal trifft, kennen wir
bereits von der Formel von Bernoulli (Bernoulli-Experiment/Binomialverteilung):

$$P(X=17) = \LoesungsRaumLang{{20 \choose 17}\cdot 0.85^{17}\cdot(1-0.85)^{20-17}}=$$
$$\LoesungsRaumLang{{20 \choose 17}\cdot 0.85^{17}\cdot (0.15)^{3}}\approx{}\LoesungsRaum{24.28\%}$$
Dabei bezeichnet $E$ das Ereignis ``siebzehn Treffer in den 20 Würfen''.

\begin{beispiel}{Kumulierte Wahrscheinlichkeit}{}
Wie groß ist nun aber die Wahrscheinlichkeit, dass er bei 20 Würfen
\textbf{höchstens} (= maximal)  17 mal trifft?
\end{beispiel}


Lösung: Wir summieren alle Wahrscheinlichkeiten auf, bei denen er
während seinen 20 Würfen
keinen Treffer, einen Treffer, zwei Treffer, ... bis zu 17 Treffer
erzielt. Dies nennen wir die kumulierte\footnote{«Kumulieren» kommt von
  lat. \texttt{cumulus} = Anhäufung.} Wahrscheinlichkeit\index{Wahrscheinlichkeit!kumulierte}:

$$P(X\le 17) = $$
\TNT{6}{$$P(X=0) + P(X=1) + P(X=2) + ... + P(X=16) + P(X=17)$$
$$={20\choose 0}\cdot 0.85^0 \cdot 0.15^{20} + {20\choose 1}\cdot 0.85^1 \cdot 0.15^{19} + {20\choose 2}\cdot 0.85^2 \cdot 0.15^{18} + ... $$

$$... + {20 \choose 17}\cdot{} 0.85^{17} \cdot{0.15}^3$$
  
}%% END TNT

Dies können wir auch mit dem Summenzeichen ($\Sigma$) schreiben:
$$P(X\le 17) = \LoesungsRaumLang{\sum_{k=0}^{17} {20 \choose k} \cdot 0.85^k \cdot 0.15^{20-k}}$$

\newpage
\textbf{Taschenrechner}:

\leserluft

Für den Taschenrechner verwenden Sie folgende Notation: 

$$\LoesungsRaumLang{\sum_{x=0}^{17}\left(( 20 \text{\,\,nCr\,\,} x ) * 0.85^x *
0.15^{20-x}\right)} \approx \LoesungsRaum{0.5951}$$




\begin{gesetz}{Kumulierte Wahrscheinlichkeit}{}
  $$P(X \le k_{\text{max}}) = \sum_{k=0}^{k_{\text{max}}} {n \choose k} \cdot{} p^k \cdot{} (1-p)^{n-k}$$


\begin{tabular}{rcl}
  $n$ &=& Anzahl Durchführungen (Baumtiefe)\\
  $p$ &=& Trefferwahrscheinlichkeit pro Durchführung\\
  $k_{\text{max}}$ &=& \textbf{Maximal} erreichte Treffer\\
\end{tabular}
\end{gesetz}


Doch auch hierzu gibt es eine eingespeicherte Formel
bei
\tiprobutton{data_stat-reg-distr} und dort unter \texttt{DISTR: 5
  Binomial\textbf{\color{red} c}df: SINGLE}

dann

$n=20$ bei $n$,

$p=0.85$ bei $p(\text{success})$ und

$k_{\text{max}}=17$ bei $x$ eintragen.

\newpage

\subsection*{Aufgaben}

\aufgabenFarbe{Wir wissen: 0.2\,\% der hergestellten Hygienemasken
  sind defekt. Wie groß ist die Wahrscheinlichkeit, dass von zufällig
  ausgewählten 1000 Stück maximal 5 defekt sind?}

\TNT{2}{$p=0.002$, $n=1000$, $P(X \le 5) \approx 98.35\,\%$}

\aufgabenFarbe{Wie groß ist die Wahrscheinlichkeit, dass
  Cristiano Ronaldo (*1985, Fußballspieler), dessen durchschnittliche Trefferwahrscheinlichkeit bei 84\% liegt,
  in einem Training mit 40 Torschüssen weniger als 36 Tore erzielt?
}
\TNT{2}{
$p=0.84$; $n=40$; $k=35$ somit $P(X < 36) = P(X\le 35) \approx 78.90\,\%$
}


\aufgabenFarbe{In einer Stadt sind 4.5\% der Autos rot. Sie stehen an der Ampel und betrachten 25 vorbeifahrende Autos.\\
  a) mit welcher Wahrscheinlichkeit ist kein rotes Auto dabei?\\
  b) mit welcher Wahrscheinlichkeit ist genau ein rotes Auto dabei?\\
  c) mit welcher Wahrscheinlichkeit sind maximal 3 rote Autos dabei (also weniger als 4 rote Autos)?
}

\TNT{4}{
    a) $31.63\%$ (binomialpdf mit $n=25; p=0.045; x=0$)\\
    b) $37.25\%$ (binomialpdf mit $n=25; p=0.045; x=1$)\\
    c) $97.57\%$ (binomial\textbf{c}df mit $n=25; p=0.045; x=3$)\\
}%% end TNT
\newpage


\aufgabenFarbe{Die Herstellung von LEDs (Licht emittierende Dioden) hat einen Ausschuss (= defekte Teile) von 0.5\%. Sie kaufen 30 solcher Dioden (LEDs) für ein Nachtlicht.
  Mit welcher Wahrscheinlichkeit sind maximal zwei davon defekt?}
\TNT{3.2}{
  $$P(X \le 2) = \sum_{k=0}^{2}{30 \choose k} \cdot{}0.005^k \cdot{} 0.995^{30-k} \approx 99.954\%$$
}%% end TNT

%%%%%%%%%%%%%%%%%%%%%%%%%%%%%%%%%%%%%%%%%%%%%%%%%%


\aufgabenFarbe{
Wie groß ist die Wahrscheinlichkeit mit einem Spielwürfel
(\epsdice{1} bis \epsdice{6}) in 10 Würfen...\\
a) ... genau viermal eine \epsdice{6} zu werfen?\\
b) ... maximal eine \epsdice{6} zu werfen?\\
c) ... nicht viermal eine eine \epsdice{6} zu werfen?\\
d) ... mehr als dreimal eine \epsdice{6} zu werfen?\\
}%% end Aufgabenfarbe
\TNTeop{
a) $P(X=4) = {10 \choose 4} \cdot{} \left( \frac16 \right)^4 \cdot{}
\left(\frac56\right)^6 \approx 5.43\%$

b) $P(X \le 1) =\sum_{k=0}^{1} {10 \choose k} \cdot{} \left( \frac16 \right)^k \cdot{}
\left(\frac56\right)^{10-k} \approx 48.45\%$

c) $P(X \ne 4) = 1- P(X = 4) \approx 1 - 0.0543 = 94.57\%$

d) $P(X > 3) = 1-P(X\le 3) = 1 - \sum_{k=0}^{3} {10 \choose k} \cdot{} \left( \frac16 \right)^k \cdot{}
\left(\frac56\right)^{10-k} \approx 6.97\%$

}%% end TNTeop (implicit \newpage)


%% Gegenwahrscheinlichkeit:
\aufgabenFarbe{Ein Frühstücksflocken-Hersteller behauptet, in jeder 10. Packung ist ein Spielzeug enthalten. Sie kaufen 10 Packungen.\\
  a) Mit welcher Wahrscheinlichkeit haben Sie mindestens zwei Spielzeuge in Ihren 10 Packungen?\\
  b) (*) Wie viele Packungen müssen Sie kaufen, damit Sie mit über 95\% Wahrscheinlichkeit mindestens zwei Spielzeuge dabei haben?}

\TNT{4}{
  Gegenwahrscheinlichkeit:
  
  a) $$P(X \ge 2) = 1 - P(X \le 1) = 1 - 73.61\% \approx 26.39\%$$
  b) 46 Stück. Pröbeln Sie mit binomialcdf mit
  $p=0.1$ und  $x=1$ solange mit $n$, bis die Wahrscheinlichkeit maximal 1 Spielzeug zu erreichen kleiner als 5\% wird. Dies ist erstmals bei 46. Packungen der Fall.
}%% end TNT



\aufgabenFarbe{Lukas bringt zu einem Würfelspiel seinen eigenen Spielwürfel mit. Im Verlauf
des Spiels entsteht der Verdacht, Lukas’ Würfel zeige überdurchschnittlich
viele Sechsen an.\\
Man beschliesst, einen Versuch durchzuführen: Es wird 30-mal gewürfelt. Es
wird vereinbart: Wenn 10-mal oder öfter eine Sechs erscheint, erklären wir
Lukas’ Würfel als gefälscht.\\
Lukas sagt: «Selbst wenn der Test zehn oder mehr Sechsen ergibt, heisst das
keineswegs, dass mein Würfel gefälscht ist; das Resultat kann auch bei einem
unverfälschten Würfel eintreten.»\\
Wie gross ist diese Wahrscheinlichkeit, dass
in 30 Würfen ein gerechter Würfel 10 oder mehr Sechsen anzeigt?
}

\TNTeop{
 $P(X \ge 10) = \sum_{k=10}^{30} {30 \choose k} \cdot{} \left( \frac16 \right)^k \cdot{}
  \left(\frac56\right)^{30-k} \approx 1.97\%$

  Mit ca. $2\%$ Wahrscheinlichkeit tritt auch bei einem gerechten Würfel ein derart
extremes Ergebnis auf. Ein gerechter Würfel wird mit ca. 98\% Wahrscheinlichkeit
ein moderateres Ergebnis zeigen (weniger als 10 Sechsen). Man wird also
eher an eine Fälschung glauben. Wer sagt, es sei ein gefälschter
Würfel irrt sich nur mit ca. $2\%$ Wahrscheinlichkeit (Hier:
\textbf{Irrtumswahrscheinlichkeit = Fehlerwahrscheinlichkeit} $\approx 2\%$).
}%% end TNTeop


\newpage
